\chapter{Suicide warning signs}
\index{Suicide warning signs}

\medskip
\noindent\textit{Educational reference; always seek professional advice for individual care.}

\section*{Definition and Overview}
Suicide warning signs are the changes in a person's words, mood, and behaviour that suggest they may be thinking about ending their life. They are signals of deep distress, and they matter because most people who die by suicide give some indication beforehand, even if those around them did not recognise it at the time. Warning signs are not always dramatic or obvious. They can be quiet and easy to dismiss, such as a tired friend who seems to be giving things away, or a colleague who has withdrawn and stopped talking about the future. Learning to recognise them, and learning what to do about them, is one of the most valuable things any person can learn.

Recognising a warning sign is only the first step. The second, and more important, step is acting on it: asking directly about suicide, staying with the person, and connecting them to help. This chapter describes the verbal, emotional, and behavioural warning signs in detail, explains the difference between warning signs and risk factors, and gives clear guidance on how to respond. It is written for anyone who may one day notice that someone they love is struggling. The chapter also carries an important reassurance: noticing the signs, asking about them, and helping someone get support can change the course of a life, and it is far better to risk a difficult conversation than to stay silent.

\section*{Epidemiology}
Around 3,000 people die by suicide in many countries each year, and for each death there are many attempts and far more people experiencing serious suicidal thoughts. Research consistently shows that a large majority of people who die by suicide communicated their distress beforehand, often to family or friends, though the messages were frequently indirect, ambiguous, or not recognised as warnings. This is why public education about warning signs is a core part of suicide prevention: the people best placed to notice the signs are the people who know the person, and most of them are not health professionals.

Warning signs are common in the weeks before a death by suicide, and they tend to cluster. Studies of people who died by suicide describe changes such as withdrawal, hopelessness, sleep disturbance, and giving away possessions in the preceding period. Men, who make up the majority of deaths, are often less likely to talk about their feelings directly and may instead show signs such as irritability, anger, or increased drinking, which are easier to miss. The good news is that awareness can be taught. Community programs that train ordinary people to recognise and respond to warning signs are associated with increased help-seeking and with lives saved, which makes this chapter directly practical.

\section*{Aetiology and Pathophysiology}
Warning signs reflect the inner state of a person in the period leading up to a suicide attempt or death. That state is usually a combination of hopelessness, the feeling of being trapped with no way out, intense psychological pain, and a growing belief that the person is a burden to others. These experiences translate into observable changes: a person loses interest in things they once cared about, withdraws from contact, stops making plans, and may begin to prepare, consciously or not, for an end to their life.

The shift from suicidal thinking to suicidal behaviour is influenced by impulsivity, intoxication, and the availability of means. Alcohol and drugs play a part in a large share of suicides by lowering inhibition and worsening mood, and acute crises often follow relationship conflict, humiliation, legal trouble, or loss. Understanding that warning signs arise from this combination of inner state and circumstance helps explain why they fluctuate, and why a person can seem to be improving just before a death: a decision to end life can bring a deceptive calm. This is why sudden improvement in someone who has been deeply distressed is itself a warning sign that should be taken seriously.

\begin{itemize}
    \item \textbf{Hopelessness:} the strongest psychological driver, expressed as a belief that nothing will get better
    \item \textbf{Psychological pain:} intense, overwhelming distress that the person feels they cannot endure
    \item \textbf{Feeling trapped:} a sense of being caught with no acceptable options
    \item \textbf{Feeling a burden:} the belief that others would be better off without them
    \item \textbf{Impulsivity and intoxication:} substances lower control and can convert thought into action
    \item \textbf{Loss and isolation:} bereavement, relationship breakdown, and disconnection sharpen risk
    \item \textbf{Access to means:} easy access to lethal methods increases the chance that an acute crisis becomes a death
\end{itemize}

\section*{Clinical Features}
Warning signs are grouped into words, mood, and behaviour. A single sign may mean little, but several signs together, especially a change from the person's usual self, demand a response.

\begin{itemize}
    \item \textbf{Talking about dying or suicide:} "I want to die", "I want to end my life", or "I wish I hadn't been born"
    \item \textbf{Talking about being a burden:} "Everyone would be better off without me"
    \item \textbf{Talking about feeling trapped:} "There's no way out" or "I can't take this anymore"
    \item \textbf{Saying goodbye:} calling or visiting to say farewell, or making unusual statements about not being around
    \item \textbf{Giving things away:} distributing possessions, money, or beloved items as if no longer needed
    \item \textbf{Withdrawal:} pulling away from family, friends, and activities, and spending more time alone
    \item \textbf{Mood changes:} deep sadness, irritability, anxiety, or rage, or a sudden unexplained calm after a period of distress
    \item \textbf{Sleep and appetite changes:} sleeping too much or too little, and eating very little
    \item \textbf{Self-neglect and risk-taking:} letting personal care slip, dangerous driving, or reckless substance use
    \item \textbf{Increased substance use:} drinking more or using drugs more heavily than usual
    \item \textbf{Seeking means:} obtaining medicines, a weapon, or information about ways to die
\end{itemize}

\section*{Risk Factors Versus Warning Signs}
It is useful to distinguish the two, because they are used in different ways.

\begin{itemize}
    \item \textbf{Risk factors:} long-standing characteristics that increase a person's overall vulnerability, such as depression, a family history of suicide, substance misuse, or a previous attempt; they tell us who is at greater risk over time but are not immediate signals
    \item \textbf{Warning signs:} recent changes in words, mood, or behaviour that suggest a person may be at risk now; they tell us when to act
    \item \textbf{Overlap:} a person with many risk factors who has recently begun talking about hopelessness and giving things away is at high risk, and warning signs in such a person are especially serious
    \item \textbf{The practical point:} risk factors help identify people who need ongoing care; warning signs identify moments when immediate action is needed
\end{itemize}

\section*{When to Seek Medical Care}
If you notice any of the warning signs described above in someone you know, act. Ask the person directly whether they are thinking about suicide, listen without judgement, and do not leave them alone if they are in danger. Then help them connect to care: a GP, a psychologist, or a crisis service. It is a myth that asking about suicide plants the idea; the opposite is true, and most people are relieved to be asked.

If you need support right now: a a a local crisis service, a local mental-health support service 1300 22 4636, a local youth crisis service 1800 55 1800 for young people, and a local culturally appropriate crisis service (13 92 76) for Aboriginal and Torres Strait Islander people.

\textbf{Contact your local emergency services for:}
\begin{itemize}
    \item A person who has just attempted suicide or harmed themselves
    \item A clear and immediate threat, with a plan and the means to carry it out
    \item Someone who has expressed suicidal intent and is now out of contact
    \item Giving away possessions or saying goodbye combined with a plan to die
\end{itemize}

\section*{Diagnosis and Investigations}
Warning signs are not a diagnosis, but recognising them leads to an assessment that is.

\begin{itemize}
    \item \textbf{Direct conversation:} asking about suicidal thoughts, plans, intent, and access to means, in a calm and caring way
    \item \textbf{Professional risk assessment:} a GP or mental health professional assesses the person's mood, thinking, and immediate risk
    \item \textbf{History:} past attempts, mental illness, substance use, and recent life events are explored
    \item \textbf{Safety planning:} if the person is at risk, a written safety plan is made, covering warning signs for the person themselves, coping strategies, contacts, and steps to safety
    \item \textbf{Removal of means:} medicines, weapons, and other lethal means are secured or removed with the person's cooperation
    \item \textbf{Follow-up:} an appointment within days, not weeks, is arranged, because the days after a crisis are the highest-risk period
    \item \textbf{Specialist referral:} referral to mental health services where risk is significant
\end{itemize}

\section*{Management}
\begin{itemize}
    \item \textbf{Take the signs seriously:} treat any mention of suicide as a real expression of distress, never as attention-seeking or a joke
    \item \textbf{Ask directly:} "Are you thinking about suicide?" asked kindly, is the single most important response to a warning sign
    \item \textbf{Listen without judgement:} allow the person to speak about what they are feeling without interrupting, arguing, or minimising
    \item \textbf{Stay with the person:} do not leave someone at risk alone while help is being arranged
    \item \textbf{Connect to professional help:} involve a GP, crisis service, or emergency department, and stay involved through the process
    \item \textbf{Secure the environment:} remove or secure medicines, alcohol, weapons, and other means
    \item \textbf{Provide a safety plan:} work with the person and a professional to write down who to call and what to do in a crisis
    \item \textbf{Follow up:} check in the next day, the next week, and keep checking; sustained connection is what makes the difference
    \item \textbf{Look after yourself:} supporting someone at risk is draining, and helpers need their own support
\end{itemize}

\section*{Complications}
\begin{itemize}
    \item \textbf{Missed opportunity:} the most important complication is a warning sign not being recognised or acted on, because it represents a lost chance to intervene
    \item \textbf{Death by suicide:} the outcome that follows when warning signs go unheeded or help arrives too late
    \item \textbf{Attempts and injury:} an attempt may occur before help is organised, with physical injury or lasting disability
    \item \textbf{Distress for the responder:} people who recognise the signs and act carry emotional weight, especially if the outcome is poor
    \item \textbf{Grief and guilt in the family:} families often replay the warning signs afterwards and may blame themselves
    \item \textbf{Contagion:} in communities and schools, a suicide can be followed by further suicidal behaviour among vulnerable people
\end{itemize}

\section*{Prognosis and Outcomes}
The prognosis depends less on the number of warning signs and more on what happens after they are recognised. When warning signs are noticed, taken seriously, and met with a direct question, professional help, and ongoing connection, the outlook is good: most people in acute suicidal crisis do not die, the crisis passes, and the underlying distress becomes treatable. People who survive an attempt and receive follow-up have a much lower risk of a further attempt, and the vast majority of those who are connected to care recover. For the person who recognises the signs, the outcome of acting is also positive: research and experience consistently show that people are grateful to be asked, and that even a difficult conversation is preferable to silence. The uncomfortable truth is that not every suicide can be prevented, but the evidence is equally clear that many, perhaps most, can be, and that the warning signs in this chapter are where prevention begins.

\section*{Prevention and Self-Care}
\begin{itemize}
    \item Learn the warning signs, and trust your instinct if you feel something is wrong with someone you know
    \item If you see the signs, ask directly about suicide; asking saves lives and cannot cause harm
    \item Know who to call: a a local crisis service, a local mental-health support service 1300 22 4636, a local youth crisis service 1800 55 1800, and a local culturally appropriate crisis service 13 92 76
    \item Keep the number of the local crisis service or GP on hand before you need it
    \item Check in on people who have recently experienced loss, a breakup, or a major setback; these are high-risk periods
    \item Reduce access to means in your home if someone is struggling, including medicines and alcohol
    \item If you have had suicidal thoughts yourself, tell someone you trust and seek help; the warning signs apply to all of us
    \item Take care of your own wellbeing while supporting others, and seek debriefing or support after a difficult intervention
\end{itemize}

\section*{Sources and Further Reading}
The following organisations provide reliable information about suicide warning signs and how to respond.

\begin{itemize}
    \item Lifeline Australia. \textit{How to recognise and respond to suicide warning signs.} https://www.lifeline.org.au/
    \item healthdirect Australia. \textit{Suicide warning signs and getting help.} https://www.healthdirect.gov.au/
    \item Beyond Blue. \textit{Supporting someone with suicidal thoughts.} https://www.beyondblue.org.au/
    \item headspace. \textit{Warning signs and support for young people.} https://headspace.org.au/
    \item Australian Department of Health and Aged Care. \textit{National suicide prevention resources.} https://www.health.gov.au/
    \item World Health Organization (WHO). \textit{Live Life: preventing suicide.} https://www.who.int/
\end{itemize}

\clearpage
